\let\negmedspace\undefined
\let\negthickspace\undefined
\documentclass[journal]{IEEEtran}
\usepackage[a5paper, margin=10mm, onecolumn]{geometry}
%\usepackage{lmodern} % Ensure lmodern is loaded for pdflatex
\usepackage{tfrupee} % Include tfrupee package

\setlength{\headheight}{1cm} % Set the height of the header box
\setlength{\headsep}{0mm}     % Set the distance between the header box and the top of the text

\usepackage{gvv-book}
\usepackage{gvv}
\usepackage{cite}
\usepackage{amsmath,amssymb,amsfonts,amsthm}
\usepackage{algorithmic}
\usepackage{graphicx}
\usepackage{textcomp}
\usepackage{xcolor}
\usepackage{txfonts}
\usepackage{listings}
\usepackage{enumitem}
\usepackage{mathtools}
\usepackage{gensymb}
\usepackage{tfrupee}
\usepackage[breaklinks=true]{hyperref}
\usepackage{tkz-euclide} 
\usepackage{listings}
% \usepackage{gvv}                                        
\def\inputGnumericTable{}                                 
\usepackage[latin1]{inputenc}                                
\usepackage{color}                                            
\usepackage{array}                                            
\usepackage{longtable}                                       
\usepackage{calc}                                             
\usepackage{multirow}                                         
\usepackage{hhline}                                           
\usepackage{ifthen}                                           
\usepackage{lscape}
\usepackage{circuitikz}
\usepackage{comment}
\tikzstyle{block} = [rectangle, draw, fill=blue!20, 
    text width=4em, text centered, rounded corners, minimum height=3em]
\tikzstyle{sum} = [draw, fill=blue!10, circle, minimum size=1cm, node distance=1.5cm]
\tikzstyle{input} = [coordinate]
\tikzstyle{output} = [coordinate]


\begin{document}

\bibliographystyle{IEEEtran}
\vspace{3cm}
\title{7.4.25}
\author{EE25BTECH11018-Darisy Sreetej}
 \maketitle
% \newpage
% \bigskip
{\let\newpage\relax\maketitle}

\renewcommand{\thefigure}{\theenumi}
\renewcommand{\thetable}{\theenumi}
\setlength{\intextsep}{10pt} % Space between text and floats


\numberwithin{equation}{enumi}
\numberwithin{figure}{enumi}
\renewcommand{\thetable}{\theenumi}

\textbf{Question}:
The locus of the centre of a circle, which touches the circle is $x^2 + y^2 -6x -6y +14 =0$ and also touches the y-axis, is given by the equation:
\begin{enumerate}
    \item $x^2-6x-10y+14=0$
    \item $x^2-10x-6y+14=0$
    \item $y^2-6x-10y+14=0$
    \item $y^2-10x-6y+14=0$
\end{enumerate}
\textbf{Solution}:\\
Given circle equation can be represented as
\begin{align}
    \norm{\vec{x}}^2+2\myvec{-3\\-3}^\top\vec{x}+14=0
\end{align}
The centre of circle is $\vec{c}_1=\myvec{3\\3}$ and radius $r=2$ $\brak{\because f= 14 ,\quad \vec{u}=\myvec{-3\\-3}} $\\
Let the centre of the moving circle be $\vec{c}=\myvec{h\\k}$\\
As the circle touches $Y-$axis , Distance of a point from $y$-axis is given by
\begin{align*}
  \text{radius }  R=h
\end{align*}
The equation of moving circle is 
\begin{align}
    \norm{\vec{x}}^2+2\myvec{-h\\-k}^\top\vec{x}+k^2=0
\end{align}
From (4.1) and (4.2) , the intersection of curves
\begin{align}
    \norm{\vec{x}}^2+2\myvec{-3\\-3}^\top\vec{x}+14 -  \norm{\vec{x}}^2-2\myvec{-h\\-k}^\top\vec{x}-k^2 =0 
    \end{align}
    \begin{align}
    2\myvec{h-3\\k-3}^\top\vec{x}= k^2-14
\end{align}
The equation of the common tangent is 
\begin{align}
\myvec{h-3\\k-3}^\top\vec{x}= \frac{k^2-14}{2}
\end{align}
Distance between the common tangent and centre of given circle equal to it's radius
\begin{align}
 2=\frac{\mydet{\myvec{h-3&k-3}\myvec{3\\3}-\frac{k^2-14}{2}}}{\sqrt{(h-3)^2 + (k-3)^2}}
\end{align}
\begin{align}
2= \frac{\mydet{3(h-3)+3(k-3)-\frac{k^2-14}{2}}}{\sqrt{(h-3)^2 + (k-3)^2}}
\end{align}
\begin{align}
2\sqrt{(h-3)^2 + (k-3)^2}= 3h +3k -\frac{k^2}{2}-11\\
\end{align}
Squaring both sides
\begin{align}
4({(h-3)^2 + (k-3)^2})= (3h +3k -\frac{k^2}{2}-11)^2
\end{align}
\begin{align}
4h^2+4k^2-24h-24k+72=\frac{k^4}{4}+9h^2+20k^2+121+18hk-66h-66k-3hk^2-3k^3
\end{align}
\begin{align}
    k^4-12k^3+(64-12h)k^2+72hk+20h^2-168h-168k+196=0
\end{align}
\begin{align}
    (k^2-10h-6k+14)(k^2-6k-2h+14)=0
\end{align}
The possible general loci are 
\begin{align}
    y^2-10x-6y+14=0 \\
    y^2-6y-2x+14=0
\end{align}
From the options , the locus of the centre of the circle is 
\begin{align}
     y^2-10x-6y+14=0
\end{align}
\begin{figure}[H]
   \centering
   \includegraphics[width=0.7\columnwidth]{figs/fig.png}
	\caption{}
   \label{}
\end{figure}
\end{document}  